\chapter[Introdução]{Introdução}\label{sec-introducao}

Com os recentes avanços tecnológicos em armazenamento e processamento de dados, tornou-se cada vez mais comum depararmo-nos com problemas em que a unidade amostral deixa de ser uma observação transversal de um sujeito e passa a ser sua trajetória ao longo do tempo. Esse paradigma é frequente em dados longitudinais e também em dados funcionais com domínio temporal \cite{diggle_analysis_2013, ramsay_functional_2005}.

Na análise de dados funcionais (FDA), cada unidade amostral corresponde a, pelo menos, uma função, possivelmente acompanhada de covariáveis (ou desfechos) escalares. Os registros costumam ser coletados com ruído e em malha discreta; por essa razão, é habitual recorrer a expansões em bases e penalizações como estratégia de pré-processamento (filtragem/suavização). A qualidade dessa etapa repercute diretamente na análise subsequente --- em especial na inferência causal funcional, em que um pré-processamento inadequado pode comprometer a identificabilidade da trajetória do efeito médio de tratamento --- ou seja, do \(\mathrm{FATE}\).

Embora teoricamente robustas, as abordagens disponíveis na literatura para estimar o \(\mathrm{FATE}\) costumam pressupor que o efeito de tratamento seja uma função suave e regular no tempo, o que suscita o uso de técnicas de pré-processamento globais, como \emph{B-Splines}. Diferentemente de outras subáreas da FDA, na inferência causal funcional as suposições de suavidade e regularidade são inverificáveis: a rigor, o efeito de tratamento pode não ser suave/regular, ainda que os desfechos observados o sejam. Somado a isso, como métodos globais tipicamente requerem acesso à curva completa para produzir estimativas estáveis, parte do efeito de tratamento pode acabar sendo atribuída a períodos anteriores à sua aplicação --- violando a ``seta do tempo''.\footnote{Em inferência causal, a premissa de que a causa precede o efeito é fundamental para a maioria das aplicações práticas. Permitir que dados do futuro influenciem representações do passado equivale, em essência, a romper essa premissa, que chamaremos de ``seta do tempo'', em referência aos grafos acíclicos direcionais (DAGs) \cite{pearlCausalDiagramsEmpirical1995}.}

Diante dessas limitações, torna-se desejável dispor de representações funcionais capazes de acomodar comportamentos localmente irregulares e, ao mesmo tempo, de reconstruir cada trecho da curva sem recorrer a observações futuras. As transformadas de ondaletas são, nesse sentido, alternativas promissoras aos métodos globais \cite{vidakovic_statistical_1999, donoho_ideal_1994}. Além de capturarem oscilações abruptas e estruturas localizadas de forma natural, elas, a depender de como implementadas, permitem organizar o processamento de modo incremental e computacionalmente eficiente, ajudando a mitigar o vazamento temporal e abrindo a possibilidade de atualizar estimativas à medida que novos dados se tornam disponíveis. Em conjunto, esses elementos sugerem um caminho promissor para sustentar análises causais funcionais em cenários mais dinâmicos e motivam as questões que orientam este trabalho.

Antes de prosseguir, convém uma nota terminológica. Ao longo deste texto, o termo ``causal'' será empregado em duas acepções: a primeira, estatística, remete ao efeito de uma intervenção sobre o desfecho (como postulado pelo \(\mathrm{FATE}\), ancorado na estrutura de resultados potenciais \cite{holland_statistics_1986, rubin_estimating_1974}); a segunda, operacional, refere-se à exigência de que os procedimentos de pré-processamento respeitem a direção temporal. Os dois usos são conceitualmente distintos, mas tornam-se intimamente relacionados quando buscamos estimar efeitos causais ao longo do tempo. Ora, se a etapa de representação funcional violar a seta do tempo, a informação vazada pode induzir dependências artificiais, comprometer a identificação e, consequentemente, a estimação do efeito causal de interesse. Desse modo, para que a causalidade no sentido estatístico seja preservada, é necessário controlar, ao menos em algum nível, o vazamento temporal --- ou seja, é preciso impor um certo nível de respeito à causalidade no sentido operacional.

Com base nessa motivação, o presente trabalho propõe o uso de transformadas de ondaletas como etapa de pré-processamento em estudos de inferência causal funcional. Mais especificamente, implementaremos ondaletas via \emph{lifting scheme} \cite{daubechies_factoring_1998, liu_improved_2014} por meio do pacote que desenvolvemos em \texttt{R}, o \texttt{rLifting} \cite{dasilvaRLiftingHighPerformanceWavelet2026}. Essa escolha torna dispensável a hipótese de suavidade do efeito e, simultaneamente, possibilita atualizar as estimativas do efeito causal em tempo quase real. Ao longo do texto, usaremos ``tempo real'' como abreviação de ``tempo quase real'', admitindo latência computacional e atualização sequencial conforme novos dados são disponibilizados.

O restante do texto está estruturado da seguinte forma. Na Seção \ref{sec-revisao}, introduzimos o estimando de interesse --- o \(\mathrm{FATE}\) (efeito médio de tratamento funcional) --- e discutimos tanto as transformadas de ondaletas em sua forma tradicional quanto sua implementação via \emph{lifting scheme}, com ênfase nas implicações para processamento incremental. A Seção \ref{sec-metodologia} detalha as adaptações necessárias em relação a trabalhos anteriores e descreve o pacote \texttt{rLifting}, que operacionaliza nossa proposta. Na Seção \ref{sec-estudos-simulacao}, especificamos o desenho dos experimentos e apresentamos os resultados da comparação com a literatura. Finalmente, encerramos com uma discussão dos achados, suas limitações e possíveis direções para trabalhos futuros.